\subsection{Global Performance Evaluator}
\label{subsec:global_evaluator}

The Global Performance Evaluator serves as the primary comparative evaluation between the \gls{snn} and \gls{cnn} architectures,
aggregating performance across the full 200-sample validation set.
Unlike the preceding targeted experiments, this benchmark establishes the \textit{Performance-Efficiency Frontier}:
quantifying the accuracy penalty incurred by transitioning to discrete spiking dynamics, whilst validating whether the multi-box refinements and others in the \texttt{ULTRA} revision of \texttt{DetectorNX}
mentioned in Section \ref{sec:architectural_evolution} are sufficient to resolve scene-level detection ambiguity.

\subsubsection{Metrics}

In this experiment, four complementary metrics are evaluated:

\begin{itemize}
    \item \textbf{\gls{miou}:} measures geometric precision between predicted and ground truth bounding boxes.

    % \item \textbf{Recall:} measures the proportion of ground truth objects successfully localised
    % at confidence threshold $\tau = 0.5$, penalising architectures that fail to respond to
    % sparse event regions %\cite{gallego2020event}.%

    \item \textbf{Recall:}
    measures the proportion of ground truth objects successfully localised.

\end{itemize}

% \textbf{Recall} measures the proportion of ground truth objects successfully localised
% at confidence threshold $\tau = 0.5$, penalising architectures that fail to respond to
% sparse event regions \parencite{gallego2020event}. \textbf{Average Detections per Frame
% (ADP)} quantifies detector selectivity relative to the ground truth mean of 2.02 objects
% per frame; significant positive deviation indicates spurious ``ghost'' detections arising
% from background noise misclassification. \textbf{Mean Prediction Confidence} captures
% network self-certainty by averaging sigmoid-activated confidence scores across all
% positive detections, probing whether spiking dynamics produce decisive or hesitant class
% activations \parencite{eshraghian2021training}.

% \subsubsection*{Hypothesis}

% We hypothesise that whilst the CNN establishes the upper mIoU bound, the SNN will
% achieve \textit{detection parity} in recall and ADP, yielding an average detection
% count consistent with the ground truth distribution. This would demonstrate that
% spiking temporal dynamics are sufficient for objectness discrimination without
% continuous-valued feature maps \parencite{roy2019towards,
% cordone2022objectdetectionspikingneural}. A measurable confidence gap is also
% anticipated: the SNN's binary spike outputs are expected to produce lower but temporally
% stable confidence distributions across the $T = 10$ timestep window, consistent with
% prior observations that surrogate-trained networks distribute representational energy
% across timesteps rather than single forward passes \parencite{fang2022resnet_sew}.

% \subsubsection*{Significance}

% This experiment provides three contributions. It delivers the \textit{parity proof},
% framing the accuracy-efficiency tradeoff quantitatively. It offers \textit{logical
% validation} that the Phase 1 ambiguity problem has been resolved through the
% architectural refinements described in Section~\ref{sec:architectural_evolution}.
% Finally, evaluation across the full validation set ensures statistical robustness,
% confirming that results are representative of the broader ETraM environment rather
% than artefacts of favourable individual frames.